\documentclass[conference]{IEEEtran}
\IEEEoverridecommandlockouts
% The preceding line is only needed to identify funding in the first footnote. If that is unneeded, please comment it out.
\usepackage{cite}
\usepackage{amsmath,amssymb,amsfonts}
\usepackage{algorithmic}
\usepackage{graphicx}
\usepackage{textcomp}
\usepackage{xcolor}

\def\BibTeX{{\rm B\kern-.05em{\sc i\kern-.025em b}\kern-.08em
    T\kern-.1667em\lower.7ex\hbox{E}\kern-.125emX}}
\begin{document}

\title{Dashcam for Traffic Object Detection\\ }

\makeatletter
\newcommand{\linebreakand}{%
  \end{@IEEEauthorhalign}
  \hfill\mbox{}\par
  \mbox{}\hfill\begin{@IEEEauthorhalign}
}
\makeatother

\author{\IEEEauthorblockN{Aura Syafa Aprilia Radim}
\IEEEauthorblockA{\textit{School of Electrical Engineering } \\
\textit{name of organization (of Aff.)}\\
City, Country \\
syafaaura@student.telkomuniversity.ac.id }
\and
\IEEEauthorblockN{Muchammad 'Irfan Chanif Rusydi}
\IEEEauthorblockA{\textit{School of Electrical Engineering } \\
\textit{name of organization (of Aff.)}\\
City, Country \\
email address or ORCID}
\and
\IEEEauthorblockN{Surya Michrandi Nasution}
\IEEEauthorblockA{\textit{dept. name of organization (of Aff.)} \\
\textit{name of organization (of Aff.)}\\
City, Country \\
email address or ORCID}
\linebreakand
% \IEEEauthorblockN{4\textsuperscript{th} Given Name Surname}
% \IEEEauthorblockA{\textit{dept. name of organization (of Aff.)} \\
% \textit{name of organization (of Aff.)}\\
% City, Country \\
% email address or ORCID}
% \and
\IEEEauthorblockN{Casi Setianingsih}
\IEEEauthorblockA{\textit{dept. name of organization (of Aff.)} \\
\textit{name of organization (of Aff.)}\\
City, Country \\
email address or ORCID}
% \and
% \IEEEauthorblockN{6\textsuperscript{th} Given Name Surname}
% \IEEEauthorblockA{\textit{dept. name of organization (of Aff.)} \\
% \textit{name of organization (of Aff.)}\\
% City, Country \\
% email address or ORCID}
}

\maketitle

\begin{abstract}
Dashcam is a camera stored in a vehicle. This tool serves to record all events in front of the vehicle. Security and safety have become a major concern in various sectors, including transportation and public security. On the highway, traffic accidents caused by the driver's ignorance of objects around the vehicle are still a serious problem. In this study, the development of a simple dashcam built from an edge computer was carried out by combining the number of cameras. Image stitching is applied to combine images that have been collected by each camera. Next, object detection is carried out on the images that have been collected. The object detection system approach is carried out using YOLOv8 which is the latest variant of the YOLO series. This research is expected to be one step in the development of an Intelligent Transportation System that is in accordance with traffic conditions in Indonesia. The results obtained in testing using the system created exist using the configuration of 78000 datasets, 3332 data validation with 8 epochs, batch size 32, linear learning rate and SGD optimization. Results are best in the morning and afternoon. The program can recognize predefined objects. 
\end{abstract}

\begin{IEEEkeywords}
object detection, YOLOv8, dashcam
\end{IEEEkeywords}

\section{Introduction}
Security and safety have become a major concern in various sectors, including transportation and public security. On the highway, traffic accidents caused by the driver's ignorance of objects around the vehicle are still a serious problem. Smart and effective object detection technology is becoming increasingly important for monitoring traffic [1]. 

Video captures from dashcams usually only show one side according to the camera position. However, this results in a lack of information about nearby objects. To overcome this problem, this study tried a solution by using two cameras or cameras placed at different positions in the vehicle. By using more than one camera, diverse viewing angles can provide more complete information about objects around the vehicle. For example, the camera on the left side helps detect objects on the left, while the camera on the right side helps detect objects on the right. This approach is expected to provide more comprehensive results according to the actual situation in front of the vehicle.
Therefore, the study of the application of dual cameras to detect objects becomes very relevant and interesting. By optimizing the use of dual cameras, it is hoped that this technology will be able to provide effective solutions in increasing driver awareness and safety on the road. The application of dual cameras for object detection has the potential to reduce accident incidents, reduce the risk of collisions, and improve the safety of all road users. 

\section{Literature Review}

\subsection{Object Detection}
Object Detection is one of the important task in computer vision field, mainly dealing with detecting instances of visual object then categorize them into several classes \cite{b2}.
With this kind of identification and localization, object detection can be used to count objects in a scene and determine and track their precise locations, all while accurately labeling them.
Object detection has been widely used for face detection, vehicle detection, pedestrian counting, web images, security systems and driverless cars.
Within the past twenty years, object detection have been going through a lot of changes and development. Although it is commonly divided into two periods: "traditional object detection" and "deep learning based". In 2012, Krizhevsky et al. \cite{b3} proposed a deep convolutional network trained on a subset of ImageNet. 

This network, called AlexNet, was the first to demonstrate that convolutional neural networks (CNNs) could be trained effectively on large-scale datasets and used to achieve state-of-the-art object detection results.
A year later, Girshick et al. proposed a new object detection framework called R-CNN \cite{b4} because it used region proposals combined with CNNs to detect objects in images.
Since then, the field of object detection has been rapidly advancing, with new models, datasets, and techniques emerging in a rapid pace. 

\subsection{YOLO (You Only Look Once)}
With the born of AlexNet, YOLO \verb|(You Only Look Once)| model was introduced in 2015. Base YOLO model can achieve 45 frame per second. While the sameller version, Fast YOLO can achieve 155 frames per second.
YOLO outperform DPM and R-CNN on Picasso Dataset and People-Art Dataset \cite{b5}.  \\

In the early period of making this paper, YOLOv7 was the latest version of YOLO. However,as January 2023, YOLOv8 was introduced by Ultralytics, the same software company that release YOLOv3 and YOLOv5. 
                                                                                                                                                                                                                                                                                                                                                                                                                                                                                                                    

\section{System Design}
Before you begin to format your paper, first write and save the content as a 
separate text file. Complete all content and organizational editing before 
formatting. Please note sections \ref{AA}--\ref{SCM} below for more information on 
proofreading, spelling and grammar.

Keep your text and graphic files separate until after the text has been 
formatted and styled. Do not number text heads---{\LaTeX} will do that 
for you.

\subsection{COCO Dataset}\label{AA}
The COCO dataset is a large-scale object detection, segmentation, and keypoint dataset. For detection dataset it consists of 80 classes, 330K images, 1.5 million object instances, 80 object categories, 91 stuff categories, 5 captions per image, 250,000 people with keypoints. 
In total, The Microsoft Common Objects in COntext contains 91 common object categories with 82 of them having more than 5,000 labeled instances. In addition, 

\subsection{Units}
\begin{itemize}
\item Use either SI (MKS) or CGS as primary units. (SI units are encouraged.) English units may be used as secondary units (in parentheses). An exception would be the use of English units as identifiers in trade, such as ``3.5-inch disk drive''.
\item Avoid combining SI and CGS units, such as current in amperes and magnetic field in oersteds. This often leads to confusion because equations do not balance dimensionally. If you must use mixed units, clearly state the units for each quantity that you use in an equation.
\item Do not mix complete spellings and abbreviations of units: ``Wb/m\textsuperscript{2}'' or ``webers per square meter'', not ``webers/m\textsuperscript{2}''. Spell out units when they appear in text: ``. . . a few henries'', not ``. . . a few H''.
\item Use a zero before decimal points: ``0.25'', not ``.25''. Use ``cm\textsuperscript{3}'', not ``cc''.)
\end{itemize}

\subsection{Equations}
Number equations consecutively. To make your 
equations more compact, you may use the solidus (~/~), the exp function, or 
appropriate exponents. Italicize Roman symbols for quantities and variables, 
but not Greek symbols. Use a long dash rather than a hyphen for a minus 
sign. Punctuate equations with commas or periods when they are part of a 
sentence, as in:
\begin{equation}
a+b=\gamma\label{eq}
\end{equation}

Be sure that the 
symbols in your equation have been defined before or immediately following 
the equation. Use ``\eqref{eq}'', not ``Eq.~\eqref{eq}'' or ``equation \eqref{eq}'', except at 
the beginning of a sentence: ``Equation \eqref{eq} is . . .''

\subsection{\LaTeX-Specific Advice}

Please use ``soft'' (e.g., \verb|\eqref{Eq}|) cross references instead
of ``hard'' references (e.g., \verb|(1)|). That will make it possible
to combine sections, add equations, or change the order of figures or
citations without having to go through the file line by line.

Please don't use the \verb|{eqnarray}| equation environment. Use
\verb|{align}| or \verb|{IEEEeqnarray}| instead. The \verb|{eqnarray}|
environment leaves unsightly spaces around relation symbols.

Please note that the \verb|{subequations}| environment in {\LaTeX}
will increment the main equation counter even when there are no
equation numbers displayed. If you forget that, you might write an
article in which the equation numbers skip from (17) to (20), causing
the copy editors to wonder if you've discovered a new method of
counting.

{\BibTeX} does not work by magic. It doesn't get the bibliographic
data from thin air but from .bib files. If you use {\BibTeX} to produce a
bibliography you must send the .bib files. 

{\LaTeX} can't read your mind. If you assign the same label to a
subsubsection and a table, you might find that Table I has been cross
referenced as Table IV-B3. 

{\LaTeX} does not have precognitive abilities. If you put a
\verb|\label| command before the command that updates the counter it's
supposed to be using, the label will pick up the last counter to be
cross referenced instead. In particular, a \verb|\label| command
should not go before the caption of a figure or a table.

Do not use \verb|\nonumber| inside the \verb|{array}| environment. It
will not stop equation numbers inside \verb|{array}| (there won't be
any anyway) and it might stop a wanted equation number in the
surrounding equation.

\subsection{Some Common Mistakes}\label{SCM}
\begin{itemize}
\item The word ``data'' is plural, not singular.
\item The subscript for the permeability of vacuum $\mu_{0}$, and other common scientific constants, is zero with subscript formatting, not a lowercase letter ``o''.
\item In American English, commas, semicolons, periods, question and exclamation marks are located within quotation marks only when a complete thought or name is cited, such as a title or full quotation. When quotation marks are used, instead of a bold or italic typeface, to highlight a word or phrase, punctuation should appear outside of the quotation marks. A parenthetical phrase or statement at the end of a sentence is punctuated outside of the closing parenthesis (like this). (A parenthetical sentence is punctuated within the parentheses.)
\item A graph within a graph is an ``inset'', not an ``insert''. The word alternatively is preferred to the word ``alternately'' (unless you really mean something that alternates).
\item Do not use the word ``essentially'' to mean ``approximately'' or ``effectively''.
\item In your paper title, if the words ``that uses'' can accurately replace the word ``using'', capitalize the ``u''; if not, keep using lower-cased.
\item Be aware of the different meanings of the homophones ``affect'' and ``effect'', ``complement'' and ``compliment'', ``discreet'' and ``discrete'', ``principal'' and ``principle''.
\item Do not confuse ``imply'' and ``infer''.
\item The prefix ``non'' is not a word; it should be joined to the word it modifies, usually without a hyphen.
\item There is no period after the ``et'' in the Latin abbreviation ``et al.''.
\item The abbreviation ``i.e.'' means ``that is'', and the abbreviation ``e.g.'' means ``for example''.
\end{itemize}
An excellent style manual for science writers is \cite{b7}.

\subsection{Authors and Affiliations}
\textbf{The class file is designed for, but not limited to, six authors.} A 
minimum of one author is required for all conference articles. Author names 
should be listed starting from left to right and then moving down to the 
next line. This is the author sequence that will be used in future citations 
and by indexing services. Names should not be listed in columns nor group by 
affiliation. Please keep your affiliations as succinct as possible (for 
example, do not differentiate among departments of the same organization).

\subsection{Identify the Headings}
Headings, or heads, are organizational devices that guide the reader through 
your paper. There are two types: component heads and text heads.

Component heads identify the different components of your paper and are not 
topically subordinate to each other. Examples include Acknowledgments and 
References and, for these, the correct style to use is ``Heading 5''. Use 
``figure caption'' for your Figure captions, and ``table head'' for your 
table title. Run-in heads, such as ``Abstract'', will require you to apply a 
style (in this case, italic) in addition to the style provided by the drop 
down menu to differentiate the head from the text.

Text heads organize the topics on a relational, hierarchical basis. For 
example, the paper title is the primary text head because all subsequent 
material relates and elaborates on this one topic. If there are two or more 
sub-topics, the next level head (uppercase Roman numerals) should be used 
and, conversely, if there are not at least two sub-topics, then no subheads 
should be introduced.

\subsection{Figures and Tables}
\paragraph{Positioning Figures and Tables} Place figures and tables at the top and 
bottom of columns. Avoid placing them in the middle of columns. Large 
figures and tables may span across both columns. Figure captions should be 
below the figures; table heads should appear above the tables. Insert 
figures and tables after they are cited in the text. Use the abbreviation 
``Fig.~\ref{fig}'', even at the beginning of a sentence.

\begin{table}[htbp]
\caption{Table Type Styles}
\begin{center}
\begin{tabular}{|c|c|c|c|}
\hline
\textbf{Table}&\multicolumn{3}{|c|}{\textbf{Table Column Head}} \\
\cline{2-4} 
\textbf{Head} & \textbf{\textit{Table column subhead}}& \textbf{\textit{Subhead}}& \textbf{\textit{Subhead}} \\
\hline
copy& More table copy$^{\mathrm{a}}$& &  \\
\hline
\multicolumn{4}{l}{$^{\mathrm{a}}$Sample of a Table footnote.}
\end{tabular}
\label{tab1}
\end{center}
\end{table}

\begin{figure}[htbp]
\centerline{\includegraphics{fig1.png}}
\caption{Example of a figure caption.}
\label{fig}
\end{figure}

Figure Labels: Use 8 point Times New Roman for Figure labels. Use words 
rather than symbols or abbreviations when writing Figure axis labels to 
avoid confusing the reader. As an example, write the quantity 
``Magnetization'', or ``Magnetization, M'', not just ``M''. If including 
units in the label, present them within parentheses. Do not label axes only 
with units. In the example, write ``Magnetization (A/m)'' or ``Magnetization 
\{A[m(1)]\}'', not just ``A/m''. Do not label axes with a ratio of 
quantities and units. For example, write ``Temperature (K)'', not 
``Temperature/K''.

\section{Results and Discussion}
\subsection[short]{Training Configuration}
All testing in this section was done in 8 epochs
\subsubitem{Learning Rate}
We are using batch 
\subsubitem{Batch Size}



\subsection{Real World Testing}

\section{Conclusion}
Dashcam has became a popular device for drivers to record the road condition. However, the dashcam video is not only used for recording the road condition, but also used for other purposes, such as the insurance claim. 
In this paper, we proposed a dashcam system to detect traffic object. The system is based on the YOLOv8 model. We also proposed a new dataset for dashcam object detection. The dataset contains 78000 images with 12 traffic objects class. 
The dataset is divided into training set and validation set. 
\section*{Acknowledgment}

The preferred spelling of the word ``acknowledgment'' in America is without 
an ``e'' after the ``g''. Avoid the stilted expression ``one of us (R. B. 
G.) thanks $\ldots$''. Instead, try ``R. B. G. thanks$\ldots$''. Put sponsor 
acknowledgments in the unnumbered footnote on the first page.

\section*{References}

Please number citations consecutively within brackets \cite{b1}. The 
sentence punctuation follows the bracket \cite{b2}. Refer simply to the reference 
number, as in \cite{b3}---do not use ``Ref. \cite{b3}'' or ``reference \cite{b3}'' except at 
the beginning of a sentence: ``Reference \cite{b3} was the first $\ldots$''

Number footnotes separately in superscripts. Place the actual footnote at 
the bottom of the column in which it was cited. Do not put footnotes in the 
abstract or reference list. Use letters for table footnotes.

Unless there are six authors or more give all authors' names; do not use 
``et al.''. Papers that have not been published, even if they have been 
submitted for publication, should be cited as ``unpublished'' \cite{b4}. Papers 
that have been accepted for publication should be cited as ``in press'' \cite{b5}. 
Capitalize only the first word in a paper title, except for proper nouns and 
element symbols.

For papers published in translation journals, please give the English 
citation first, followed by the original foreign-language citation \cite{b6}.

\begin{thebibliography}{00}
\bibitem{b1} G. Eason, B. Noble, and I. N. Sneddon, ``On certain integrals of Lipschitz-Hankel type involving products of Bessel functions,'' Phil. Trans. Roy. Soc. London, vol. A247, pp. 529--551, April 1955.
\bibitem{b2} https://arxiv.org/pdf/1905.05055.pdf
\bibitem{b3} Krizhevsky, A., Sutskever, I., Hinton, G. E. (2012). ImageNet classification with deep convolutional neural networks. Communications of the ACM, 60(6), 84–90. https://doi.org/10.1145/3065386
\bibitem{b4} https://arxiv.org/abs/1311.2524
\bibitem{b5} https://arxiv.org/abs/1506.02640
\bibitem{b6} R. Nicole, ``Title of paper with only first word capitalized,'' J. Name Stand. Abbrev., in press.
\bibitem{b7} Y. Yorozu, M. Hirano, K. Oka, and Y. Tagawa, ``Electron spectroscopy studies on magneto-optical media and plastic substrate interface,'' IEEE Transl. J. Magn. Japan, vol. 2, pp. 740--741, August 1987 [Digests 9th Annual Conf. Magnetics Japan, p. 301, 1982].
\bibitem{b8} M. Young, The Technical Writer's Handbook. Mill Valley, CA: University Science, 1989.
\end{thebibliography}
\vspace{12pt}
\color{red}
IEEE conference templates contain guidance text for composing and formatting conference papers. Please ensure that all template text is removed from your conference paper prior to submission to the conference. Failure to remove the template text from your paper may result in your paper not being published.

\end{document}
